\subsection{基本域与稳定子群}

回顾《积分》第 VII 章 §2 第 10 号定义 2：若 E 的子集 D 使 W 在 E 中的每个轨道都恰好与 D 相交于一点，则称 D 是群 W 的基本域。这等价于下列两个条件：

\begin{enumerate}
\def\labelenumi{\alph{enumi})}
\item
对于每个 \(x \in \mathrm{E}\)，存在 \(w \in \mathrm{W}\)，使得 \(w(x) \in \mathrm{D}\)。
\item
若 \(x, y \in D\) 且 \(w \in W\) 满足 \(y = w(x)\)，则 y = x（即使 \(w \neq 1\)）。
\end{enumerate}

我们将证明下列三个命题：

定理 2. 对任意室 C，C 的闭包 \(\overline{\mathbf{C}}\) 是 W 在 E 上作用的一个基本域。

命题 1. 设 \(\mathbf{F}\) 是一个面片，\(\mathbf{C}\) 是满足 \(\mathbf{F} \subset \overline{\mathbf{C}}\) 的室。设 \(w \in W\)。下列条件等价：

\begin{enumerate}
\def\labelenumi{(\roman{enumi})}
\item
\(w(\mathbf{F})\) 与 F 相交；
\item
\(w(\mathbf{F}) = \mathbf{F};\)
\item
\(w(\overline{\mathrm{F}}) = \overline{\mathrm{F}};\)
\item
w 至少固定 F 中的一个点；
\item
w 固定 F 的每一个点；
\item
w 固定 \(\overline{\mathbf{F}}\) 的每一个点；
\item
w 属于 W 中由包含 F 的 C 的墙的反射生成的子群。
\end{enumerate}

对于 E 的每个子集 A，记 \(W(A)\) 为 W 中由固定 A 的每一个点的元素所成的子群。

命题 2. 设 A 是 E 的非空子集，令 \(\mathfrak{H}_{\mathrm{A}}\) 为包含 A 的超平面 \(\mathrm{H} \in \mathfrak{H}\) 所成的集合，令 \(\mathrm{A}'\) 为这些 \(\mathrm{H} \in \mathfrak{H}_{\mathrm{A}}\) 的交，并令 F 为 E 中在 \(\mathrm{A}'\) 内部开着的一个面片（§1，第 3 号，命题 7）。则 \(\mathrm{W(A)} = \mathrm{W(A')} = \mathrm{W(F)}\)，且 \(\mathrm{W(A)}\) 由关于 \(\mathfrak{H}_{\mathrm{A}}\) 中超平面的反射生成。

我们先证明下列断言：

\begin{enumerate}
\def\labelenumi{(\Roman{enumi})}
\tightlist
\item
设 C 是一个室，设 x 和 y 是 \(\overline{C}\) 中的两个点，并设 \(w \in W\) 满足 \(w(x) = y\)。则 x = y，且 w 属于子群 \(W_{N}\)，其中 N 是包含 x 的 C 的墙的集合。
\end{enumerate}

我们对长度 \(q\) 的 \(w\)（相对于集合 \(S\)（由关于 \(C\) 的墙的反射组成））作归纳，\(q = 0\) 时结论显然。若 \(q \geqslant 1\)，则存在 \(H\) 为 \(C\) 的一个墙以及元素 \(w' \in W\)，使得 \(w = s_H w'\) 且 \(l(w') = q - 1\)。由于 \(l(s_H w) < l(w)\)，根据第 2 号的定理 1，室 \(C\) 与 \(w(C)\) 位于 \(H\) 的相反两侧。因此，\(\overline{C} \cap w(\overline{C}) \subset H\)，所以 \(y \in H\)。于是 \(y = w'(x)\)，归纳假设推出 \(x = y\) 且 \(w' \in W_{\mathfrak{N}}\)。由于 \(y \in H\)，可知 \(H \in \mathfrak{N}\)，从而 \(w = s_H w' \in W_{\mathfrak{N}}\)，(I) 的证明完成。

定理 2 的证明：这由 (I) 和引理 2 得出。

命题 1 的证明：我们知道两个不同的面片不相交且闭包不同（§1，第 2 号，命题 3 的推论）。因此 (i)、(ii)、(iii) 等价。另一方面，显然

\[
(\mathrm{vii}) \Longrightarrow (\mathrm{vi}) \Longrightarrow (\mathrm{v}) \Longrightarrow (\mathrm{iv}) \Longrightarrow (\mathrm{i})
\]

且断言 (I) 表明 (i) \(\Longrightarrow\) (vii)。

命题 2 的证明：令 \(\mathbf{A}''\) 为由 A 生成的 \(\mathbf{E}\) 的仿射子空间。显然，\(\mathrm{W(A)} = \mathrm{W(A'')}\)。由 §1 第 3 号的命题 7，存在点 \(x \in \mathbf{A}''\)，它不属于任何超平面 \(\mathbf{H} \in \mathfrak{H} - \mathfrak{H}_{\mathbf{A}}\)。令 \(\mathbf{F}_x\) 为包含 \(x\) 的面片：它在 \(\mathbf{A}''\) 中是开的，而命题 1 表明

\[
\mathrm{W} \left(\mathrm{F} _ {x}\right) \subset \mathrm{W} \left(\mathrm{A} ^ {\prime}\right) \subset \mathrm{W} (\mathrm{A}) = \mathrm{W} \left(\mathrm{A} ^ {\prime \prime}\right) \subset \mathrm{W} (\{x \}) = \mathrm{W} \left(\mathrm{F} _ {x}\right),
\]

从而 \(\mathrm{W(A) = W(A') = W(F_x)}\)。将 A 替换为 F，我们也有

\[
\mathrm{W} (\mathrm{A}) = \mathrm{W} (\mathrm{F}),
\]

因此命题得证。

注. 1) 由定理 2 可知，若 C 是一个室而 F 是一个面片，则存在唯一的面片 \(F'\) 包含于 \(\overline{C}\)，它被 W 的一个元素变换为 F。

\begin{enumerate}
\def\labelenumi{\arabic{enumi})}
\setcounter{enumi}{1}
\item
由命题 1 和 2 可知，对于 E 的任意非空子集 A，存在一点 \(a \in E\)，使得 \(W(A) = W(\{a\})\)；此外，群 \(W(A)\) 是一个 Coxeter 群（定理 1）。
\item
设 \(\mathbf{C}\) 是 \(\mathbf{E}\) 的一个室，\(\mathbf{S}\) 是关于 \(\mathbf{C}\) 的墙的反射所成的集合。设 \(w \in W\)，并设 \((s_1, \ldots, s_q)\) 是 \(w\) 关于 \(\mathbf{S}\) 的一个约化分解。若 \(x \in \overline{\mathbf{C}}\) 被 \(w\) 固定，则有 \(s_j(x) = x\) 对所有 \(j\) 成立：这由上述命题 1 以及《代数》第 IV 章 §1 的命题 7 的推论 1 得出。
\end{enumerate}

